
on ait : \(a_i^{(m)}=P_i^{(m,r)}(a_r,\ldots,a_i)\).

\textbf{Lemme 2.11.} \emph{Soient \(S\) un schéma noethérien et \(X\) un \(S\) schéma projectif, muni d'un faisceau très ample \(H={O}_X(1)\). On suppose que les fibres de \(X/S\) sont géométriquement intègres, de même dimension \(r\). Soit \(\mathcal L\) une famille de classes de faisceaux inversibles \(L_K\) sur les fibres de \(X/S\). Alors pour que \(\mathcal L\) soit limitée, (il faut et) il suffit que dans les polynômes de Hilbert \(\chi(L_K(n))=\sum_{i=0}^{r}a_i\binom{n+i}{i}\) des \(L_K\in\mathcal L\), tous les coefficients \(a_i\) restent bornés. De plus, si les fibres de \(X/S\) sont aussi géométriquement normales, alors il suffit que les trois coefficients \(a_r,a_{r-1},a_{r-2}\) restent bornés.}

En effet, soit \(\delta=\operatorname{supp}\{d(L_K)\}\). Si \(\delta>0\), on peut remplacer \(\mathcal L\) par la famille des classes des \(L_K(-\delta)\) grâce à 2.3 et 2.10\,; on peut ainsi supposer \(\delta\le0\). Alors \(\mathcal L\) est une \((b)\)-famille d'après 2.5, donc limitée si tous les coefficients \(a_i\) restent bornés, en vertu de 1.13, et, dans le cas où les fibres de \(X/S\) sont géométriquement normales, il suffit pour cela que les trois coefficients \(a_r,a_{r-1},a_{r-2}\) restent bornés, en vertu de 2.9.

\subsection*{3. Théorèmes de finitude généraux}
\textbf{Lemme 3.1.} \emph{Soient \(S\) un sch\'ema et \(X\) un \(S\)-sch\'ema quasi-compact et quasi-s\'epar\'e\,; soit \(\underline{\Pic}_{X/S}\) le foncteur de Picard, (localis\'e pour la topologie f.p.p.f.)}\SGAAuthorFootnote{*}{Cf. SGA 3 IV 6.3, et FGA n\textsuperscript{o} 232 pour la d\'efinition de \(\underline{\Pic}_{X/S}\), o\`u il convient ici, comme indiqu\'e, de remplacer la topologie fpqc par la topologie fppf.}.
\emph{Alors pour tout syst\`eme projectif filtrant \((Z_\alpha)\) de \(S\)-sch\'emas 
